% บทคัดย่อภาษาไทย
% หมายเหตุ: ไม่ควรเกิน 300 คำ
% ควรครอบคลุม: วัตถุประสงค์ → วิธีการ → ผลลัพธ์ → การประยุกต์ใช้
%
โครงงานนี้มีวัตถุประสงค์เพื่อออกแบบและพัฒนาแอปพลิเคชันบนมือถือชื่อ ``Pacegasus'' สำหรับส่งเสริมพฤติกรรมการวิ่งออกกำลังกายอย่างยั่งยืน โดยผสานแนวคิดเกมมิฟิเคชัน (Gamification) เข้ากับการปรับแผนฝึกอัตโนมัติและระบบปฏิสัมพันธ์ทางสังคม บนพื้นฐานทฤษฎีการกำหนดใจตนเอง (Self-Determination Theory: SDT) และเทคนิคการเปลี่ยนพฤติกรรม (Behavior Change Techniques: BCTs)

ระบบพัฒนาบนแพลตฟอร์ม React Native รองรับระบบปฏิบัติการ Android โดยใช้เซนเซอร์ GPS และ Accelerometer ในตัวสมาร์ตโฟนเพื่อติดตามระยะทาง ความเร็วเฉลี่ย (Pace) และจำนวนก้าว โดยไม่จำเป็นต้องพึ่งพาอุปกรณ์เสริมภายนอก ฝั่งเซิร์ฟเวอร์ใช้ Node.js และ PostgreSQL สำหรับประมวลผลและจัดเก็บข้อมูล

ฟีเจอร์หลักของระบบประกอบด้วยเอนจินปรับแผนฝึกอัตโนมัติ (Adaptive Training Engine) ที่วิเคราะห์ข้อมูลสุขภาวะรายวัน ค่า Session-RPE และอัตราส่วนภาระฝึกระยะเฉียบพลันต่อเรื้อรัง (ACWR) เพื่อปรับมิชชั่นรายวันและรายสัปดาห์ให้เหมาะสมกับแต่ละบุคคล ระบบเกมมิฟิเคชันที่ประกอบด้วยคะแนน เหรียญตรา กระดานจัดอันดับ และโหมดเนื้อเรื่อง ตลอดจนระบบสังคมที่เปิดให้ผู้ใช้เพิ่มเพื่อน เข้าร่วมกิจกรรมกลุ่ม และแชร์ผลการวิ่ง

ผลที่คาดว่าจะได้รับจากโครงงานนี้คือต้นแบบแอปพลิเคชันที่ใช้งานได้จริง ซึ่งสามารถช่วยเพิ่มความสม่ำเสมอในการออกกำลังกาย ลดความเสี่ยงจากการบาดเจ็บ และส่งเสริมพฤติกรรมสุขภาพที่ยั่งยืนในระยะยาว อันเป็นแนวทางหนึ่งในการบรรเทาปัญหาโรคไม่ติดต่อเรื้อรัง (NCDs) ในระดับประชากร
